% GENERATED FILE. Edit canonical lessons, then run npm run generate:books.
% canonical-source-hash: fnv1a64:6528e85493aa4653
% canonical-lessons: ZH-C219-yun, ZH-W219-yun-cloud, ZH-C219-duoyun, ZH-C219-jiangwen, ZH-C219-qiwen, ZH-C219-fangqing

\chapter{Cloudy Tomorrow}
\label{ch:zh-cloudy-tomorrow}
\hlchaptermodality{219}

\begin{chapteropening}
\textbf{By the end of this chapter:} \emph{I can say there will be cloud, a cloudy day, a cold snap, what the temperature is and when it will clear.}

Complete the last lesson of chapter 219: I can say there will be cloud, a cloudy day, a cold snap, what the temperature is and when it will clear.
\end{chapteropening}

\section[\texorpdfstring{\zh{云} (yún)}{云 (yún)}]{\zh{云} (yún) — cloud}

\label{lesson:ZH-C219-yun}



\begin{quote}
Before the new one: say the Mandarin for dry air, then the Mandarin for damp.
\end{quote}

\subsection*{You'll want to know: \zh{云}}
\textbf{\zh{云}} — \emph{yún} — \textquotedblleft{}cloud\textquotedblright{}.

\begin{quote}
\textbf{\zh{明天下午会有云}} \emph{míngtiān xiàwǔ huì yǒu yún} — There'll be some cloud tomorrow afternoon.
\end{quote}

The new character, in pieces: \zh{云} — four strokes: two bars, then a stroke that falls to the left and turns right along the foot, and a dot. The next lesson puts it in your hand.

\begin{cousinweb}
\textbf{\zh{云}} (\emph{yún}) is inherited from Middle Chinese \emph{hjun} (Guangyun, in Baxter's transcription), which explains it, after an older book, as the vapour of mountains and rivers. The dictionary writes the word for cloud with \zh{雨} on top, and lists \zh{云} as the older way of writing it.
\end{cousinweb}

\begin{grammarlens}[title={What you've built}]
One more word for talking about plans.
\end{grammarlens}

\subsection*{Your turn}
\begin{itemize}
\raggedright
  \item \emph{Say it:} \emph{yún}
  \item \emph{Say it:} \emph{yún}, once more
  \item \emph{Say it:} \emph{míngtiān xiàwǔ huì yǒu yún} — There'll be some cloud tomorrow afternoon.
  \item \emph{Recall:} say the Mandarin for dry air, then the Mandarin for damp, then say \emph{yún} again
\end{itemize}

\begin{tcolorbox}[breakable,colback=teal!4,colframe=teal!35!black,title={Before you move on}]
What does \zh{云} mean? (\textquotedblleft{}Cloud\textquotedblright{}.) Say it once more.
\end{tcolorbox}
\section[\texorpdfstring{\zh{云} (yún)}{云 (yún)}]{\zh{云} — two bars, then a turn and a dot}

\label{lesson:ZH-W219-yun-cloud}



\begin{quote}
The word first, then the character inside it.

\begin{itemize}
\raggedright
  \item \emph{Recall:} say \emph{yún} — \textbf{R1}, one lesson back, before it decays
\end{itemize}

\begin{quote}
\zh{二} \zh{雨}
\end{quote}

\emph{Read it:} each one aloud, with nothing said first — \emph{èr}, \emph{yǔ}
\end{quote}

\subsection*{Script — two bars, then a turn and a dot}
\hlblockfigure{\detokenize{figures/ZH-W219-yun-cloud-filmstrip.pdf}}{How \zh{云} is written, stroke by stroke}

\begin{quote}
\zh{云}
\end{quote}

\textbf{Two pieces, top before bottom.}

\textbf{On top, \zh{二}, two strokes:} a short bar, then a longer one under it.

\textbf{Under it, two strokes.} A stroke falling to the left that turns and runs right along the foot. Then a dot at the right, above the end of that stroke, falling to the right.

2 + 2 = four.

\subsection*{Your turn}
\begin{itemize}
\raggedright
  \item \emph{Write it:} \zh{云}, all four, with the model in front of you
  \item \emph{Cover:} write it again from memory, and count the pieces, not the strokes
\end{itemize}

\begin{tcolorbox}[breakable,colback=teal!4,colframe=teal!35!black,title={Before you move on}]
How many strokes in \zh{云}? (\textbf{Four}.) What is on top? (\textbf{\zh{二}, two bars}.) What does the third stroke do? (\textbf{It falls to the left and turns right}.)

Source: \href{https://raw.githubusercontent.com/chanind/hanzi-writer-data/68d10a4b21150cae5e1ebbd223eed289cf32d90c/data/\%E4\%BA\%91.json}{Hanzi Writer Data \zh{云}, pinned stroke paths and medians}.
\end{tcolorbox}
\section[\texorpdfstring{\zh{多云} (duōyún)}{多云 (duōyún)}]{\zh{多云} (duōyún) — cloudy}

\label{lesson:ZH-C219-duoyun}



\begin{quote}
Before the new one: say the Mandarin for damp, then the Mandarin for a cloud.
\end{quote}

\subsection*{You'll want to know: \zh{多云}}
\textbf{\zh{多云}} — \emph{duōyún} — \textquotedblleft{}cloudy\textquotedblright{}.

\begin{quote}
\textbf{\zh{明天多云风不大}} \emph{míngtiān duōyún fēng bú dà} — Tomorrow will be cloudy, without much wind.
\end{quote}

\begin{cousinweb}
\textbf{\zh{多云}} (\emph{duōyún}) is built from \zh{多} 'many' and \zh{云} 'cloud', both inherited from Middle Chinese, \emph{ta} and \emph{hjun} (Guangyun, in Baxter's transcription).
\end{cousinweb}

\begin{grammarlens}[title={What you've built}]
One more word for talking about plans.
\end{grammarlens}

\subsection*{Your turn}
\begin{itemize}
\raggedright
  \item \emph{Say it:} \emph{duōyún}
  \item \emph{Say it:} \emph{duōyún}, once more
  \item \emph{Say it:} \emph{míngtiān duōyún fēng bú dà} — Tomorrow will be cloudy, without much wind.
  \item \emph{Recall:} say the Mandarin for damp, then the Mandarin for a cloud, then say \emph{duōyún} again
\end{itemize}

\begin{tcolorbox}[breakable,colback=teal!4,colframe=teal!35!black,title={Before you move on}]
What does \zh{多云} mean? (\textquotedblleft{}Cloudy\textquotedblright{}.) Say it once more.
\end{tcolorbox}
\section[jiàngwēn]{jiàngwēn — to turn colder; to cool down}

\label{lesson:ZH-C219-jiangwen}



\begin{quote}
Before the new one: say the Mandarin for a cloud, then the Mandarin for cloudy.
\end{quote}

\subsection*{You'll want to know: jiàngwēn}
\emph{jiàngwēn} — \textquotedblleft{}to turn colder; to cool down\textquotedblright{}.

\begin{quote}
\emph{hòutiān huì jiàngwēn} — It's going to turn colder the day after tomorrow.
\end{quote}

Its characters come later; for now it is a word for the ear.

\begin{cousinweb}
\emph{jiàngwēn} is built from \emph{jiàng} 'to come down' and \emph{wēn} 'warm', both inherited from Middle Chinese, \emph{kaewngH} and \emph{'won} (Guangyun, in Baxter's transcription), which glosses the first 'down' and 'to fall', and the second, among other senses, 'mild' and 'warm'.
\end{cousinweb}

\begin{grammarlens}[title={What you've built}]
One more word for talking about plans.
\end{grammarlens}

\subsection*{Your turn}
\begin{itemize}
\raggedright
  \item \emph{Say it:} \emph{jiàngwēn}
  \item \emph{Say it:} \emph{jiàngwēn}, once more
  \item \emph{Say it:} \emph{hòutiān huì jiàngwēn} — It's going to turn colder the day after tomorrow.
  \item \emph{Recall:} say the Mandarin for a cloud, then the Mandarin for cloudy, then say \emph{jiàngwēn} again
\end{itemize}

\begin{tcolorbox}[breakable,colback=teal!4,colframe=teal!35!black,title={Before you move on}]
What does jiàngwēn mean? (\textquotedblleft{}To turn colder; to cool down\textquotedblright{}.) Say it once more.
\end{tcolorbox}
\section[qìwēn]{qìwēn — air temperature}

\label{lesson:ZH-C219-qiwen}



\begin{quote}
Before the new one: say the Mandarin for cloudy, then the Mandarin for to turn colder.
\end{quote}

\subsection*{You'll want to know: qìwēn}
\emph{qìwēn} — \textquotedblleft{}air temperature\textquotedblright{}.

\begin{quote}
\emph{jīntiān de qìwēn bǐ zuótiān gāo} — It's warmer today than yesterday.
\end{quote}

Its characters come later; for now it is a word for the ear.

\begin{cousinweb}
\emph{qìwēn} is built from \zh{气} 'air' and \emph{wēn} 'warm', both inherited from Middle Chinese, \emph{khj+jH} and \emph{'won} (Guangyun, in Baxter's transcription); the dictionary glosses the first 'breath'.
\end{cousinweb}

\begin{grammarlens}[title={What you've built}]
One more word for talking about plans.
\end{grammarlens}

\subsection*{Your turn}
\begin{itemize}
\raggedright
  \item \emph{Say it:} \emph{qìwēn}
  \item \emph{Say it:} \emph{qìwēn}, once more
  \item \emph{Say it:} \emph{jīntiān de qìwēn bǐ zuótiān gāo} — It's warmer today than yesterday.
  \item \emph{Recall:} say the Mandarin for cloudy, then the Mandarin for to turn colder, then say \emph{qìwēn} again
\end{itemize}

\begin{tcolorbox}[breakable,colback=teal!4,colframe=teal!35!black,title={Before you move on}]
What does qìwēn mean? (\textquotedblleft{}Air temperature\textquotedblright{}.) Say it once more.
\end{tcolorbox}
\section[fàngqíng]{fàngqíng — to clear up (of the sky)}

\label{lesson:ZH-C219-fangqing}



\begin{quote}
Before the new one: say the Mandarin for to turn colder, then the Mandarin for air temperature.
\end{quote}

\subsection*{You'll want to know: fàngqíng}
\emph{fàngqíng} — \textquotedblleft{}to clear up (of the sky)\textquotedblright{}.

\begin{quote}
\emph{xiàwǔ huì fàngqíng} — It'll clear up in the afternoon.
\end{quote}

Its characters come later; for now it is a word for the ear.

\begin{cousinweb}
\emph{fàngqíng} is built from \emph{fàng} 'to let go' and \emph{qíng} 'clear'. The first is inherited from Middle Chinese \emph{pjangH} (Guangyun, in Baxter's transcription), glossed 'to drive out'; the dictionary glosses the second, \emph{dzjeng} there, 'a clear sky'.
\end{cousinweb}

\begin{grammarlens}[title={What you've built}]
One more word for talking about plans.
\end{grammarlens}

\subsection*{Your turn}
\begin{itemize}
\raggedright
  \item \emph{Say it:} \emph{fàngqíng}
  \item \emph{Say it:} \emph{fàngqíng}, once more
  \item \emph{Say it:} \emph{xiàwǔ huì fàngqíng} — It'll clear up in the afternoon.
  \item \emph{Recall:} say the Mandarin for to turn colder, then the Mandarin for air temperature, then say \emph{fàngqíng} again
\end{itemize}

\begin{tcolorbox}[breakable,colback=teal!4,colframe=teal!35!black,title={Before you move on}]
What does fàngqíng mean? (\textquotedblleft{}To clear up (of the sky)\textquotedblright{}.) Say it once more.
\end{tcolorbox}
